\chapter{Parallelisation design}
\label{ch:parallelisation}

\section{One interface}

Every backend implements the same small interface, and it is deliberately the
minimal common denominator of the models it spans: a data parallel map, a
reduction, a barrier, and the few extra operations a distributed memory model
genuinely needs. Anything richer would privilege one model and stop the
comparison being fair. The shape follows the structured parallel patterns
treatment of McCool, Robison and Reinders \cite[Ch.~3, 5]{mccool2012}.

\begin{lstlisting}[caption={The parallelism interface, abbreviated}]
void parallel_for(Index n, const RangeBody& body);
Real reduce(Index n, Real init, const RangeReducer& reducer);
void barrier();

Range local_rows(Index total_rows) const;
void exchange_halo(VectorView grid, Index stride, Index rows);
void gather_rows(VectorView data, Range local);
void run_ordered(const std::function<void()>& work, bool forward, ...);
\end{lstlisting}

\subsection{Why the callbacks are chunk level}

A body receives a \texttt{Range} and loops over it itself, rather than being
called once per element. Two consequences follow, both deliberate.

First, the \texttt{std::function} indirection is paid once per chunk rather than
once per element, so it is $O(\text{workers})$ per sweep and does not
contaminate the timings the study exists to produce. Second, the innermost loop
stays a plain loop over contiguous memory that the compiler can still vectorise,
so each backend measures its own dispatch and synchronisation cost rather than a
penalty the abstraction imposed on all of them equally.

\subsection{What backends may not do}

No backend leaks its model's types through the interface: there is no
\texttt{MPI\_Comm}, no \texttt{omp\_} type and no \texttt{cudaStream\_t} in any
signature. The CUDA path goes further and crosses a C ABI boundary, which it has
to for reasons given in Section \ref{sec:cuda-design}.

\section{Determinism as a design decision}
\label{sec:determinism}

Floating point addition is not associative, so the result of a reduction depends
on the order in which the partial sums are combined. Most parallel libraries
treat this as an unavoidable nuisance and paper over it with a tolerance. This
library treats it as a choice.

Under \texttt{ReductionMode::Deterministic}, the chunk grid used by every
reduction is fixed by the problem size alone:

\begin{lstlisting}[caption={The chunking policy that makes determinism possible}]
inline constexpr Index DETERMINISTIC_CHUNKS = 512;

constexpr Index reduction_chunk_count(Index n) noexcept {
    return n <= 0 ? 0 : std::min(DETERMINISTIC_CHUNKS, n);
}
\end{lstlisting}

Each chunk's partial lands in its own slot and the slots are summed in ascending
index order afterwards. Which worker computed which slot, and in what order they
finished, therefore cannot affect the answer. The consequence is that a
reduction computed on one thread, on twenty threads, or on a persistent pool
returns bit identical results, and the equivalence suite can assert exact
equality rather than hiding a real disagreement under a tolerance.

Two limits are documented rather than concealed. Across MPI rank counts the
results agree to reduction tolerance and not bitwise, because rank boundaries
are chosen for load balance and therefore do not align with the fixed chunk
grid, which regroups the additions; within one rank count an ordered allgather
makes the result reproducible and identical on every rank. On the GPU the
reductions use a tree inside each thread block, which is not an ordered sum and
cannot be made into one by a compiler flag; GPU \emph{sweeps} are bit identical
to the CPU, and anything passing through a reduction is compared to tolerance.

The cost of this choice is measured, not waved away: Section
\ref{sec:reduction-cost} prices the deterministic mode against each model's
native reduction.

\section{Contraction, and a claim that nearly was not true}

Getting the GPU sweeps to agree bit for bit with the CPU required disabling
fused multiply add contraction on both sides: \verb|--fmad=false| for
\texttt{nvcc} and \texttt{-ffp-contract=off} for the host compiler. (Those flags
are set in \verb|\verb| spans throughout, because a double hyphen in ordinary
prose typesets as exactly the character ground rule 1 forbids.)

The reason is instructive. The Jacobi update has the form
$\tfrac{1}{4}(b + \sum \text{neighbours})$, which offers no opportunity to fuse,
and it agreed immediately. The SOR update has the form $ab + cd$, which a
compiler may legally contract into a fused multiply add, and the host was
contracting it while the device was not. The disagreement was one bit, and it
would have been easy to absorb into a tolerance and never think about again.

It was not absorbed, because a claim of identical numerics that holds only when
the compiler happens not to contract is a much weaker claim than it appears, and
one that would quietly break on a compiler upgrade. On kernels that are
bandwidth bound the cost of disabling contraction is not measurable.

\section{What each backend does underneath}

\subsection{Shared memory}

\textbf{OpenMP} is the only backend that does not partition the range itself. It
hands the chunk grid to OpenMP's own worksharing construct, so what is measured
is OpenMP's scheduler and barrier rather than a hand written partitioner wearing
an OpenMP hat. The build asserts the \texttt{\_OPENMP} version macro at compile
time. The reference document cited is OpenMP 6.0 \cite{openmp2024}; GCC 16
reports 202111, which is OpenMP 5.2, and the code is deliberately restricted to
what the compiler actually implements.

\textbf{POSIX threads} uses raw \texttt{pthread\_create}, a mutex and condition
variable for dispatch, a counter and a second condition variable for completion,
and \texttt{pthread\_setaffinity\_np} for binding. The pool is persistent:
creating threads per sweep would measure \texttt{pthread\_create}, which is a
well known number and not the one this study is about. Workers wait on a
generation counter rather than a flag, so a task published while a worker was
still finishing the previous one cannot be missed, which is the lost wakeup race
a plain boolean would have.

\textbf{The jthread pool} uses only what C++23 provides: \texttt{std::jthread}
workers, \texttt{std::stop\_token} shutdown, and a pair of \texttt{std::barrier}
phases entered in strict alternation. Using one barrier for both release and
collect would let a fast worker race into the next task before a slow one had
left the previous. Work stealing is deliberately absent: the loops here are
uniform stencil sweeps over contiguous memory, where a static partition is
already balanced and a stealing deque would add per task atomics that buy
nothing. Section \ref{sec:schedule-cost} provides the measurement behind that
choice rather than the assertion of it.

\subsection{Distributed memory}

The MPI backend uses remainder aware block row decomposition, so block sizes
differ by at most one and no row is dropped. The halo exchange is two
\texttt{MPI\_Sendrecv} calls moving one grid row each; every MPI call goes
through a checking macro, so a failure becomes a diagnosed exception carrying
its call site rather than a silently wrong answer.

Each rank allocates the whole grid and computes only its own band. A production
code would allocate the local slab only, and this one would too if memory were
the constraint; at 4096 squared a vector is 134 MB, which fits comfortably. What
matters for the study is that the communication volume is identical either way,
and the choice keeps the solver code identical across every backend.

The installed OpenMPI 5.0.10 advertises MPI 3.1 through \texttt{MPI\_VERSION}
while implementing many MPI 4.0 entry points \cite{mpi2021}. The build probes for
symbols rather than testing the version macro, since guarding on the macro would
disable functionality that demonstrably works.

The hybrid backend \emph{inherits} from the MPI backend rather than
reimplementing it, so the halo exchange, the ordered pass and the deterministic
reduction are literally the same code. Any difference the sweep measures between
the two is therefore the threading and nothing else.

\subsection{Sequential dependence across ranks}

Natural ordering Gauss Seidel and SOR are sequentially dependent by definition.
Preserving that across a row decomposition means rank $r$ may not start until
rank $r-1$ has finished and handed over its boundary row, which is the classical
pipelined Gauss Seidel. The library implements exactly that, through
\texttt{run\_ordered}.

The alternative would have been to substitute a red black reordering silently
whenever more than one worker was present. That would have produced attractive
scaling curves for a method that does not scale, by changing the method. Doing
it honestly means these solvers are bit identical to the serial backend at every
rank count \emph{and} show essentially no speedup, and the report reports both.

\section{The CUDA boundary}
\label{sec:cuda-design}

CUDA 13.3 refuses host compilers newer than GCC 15, and cannot parse GCC 15's
own libstdc++ headers either, because they use the C++23 \texttt{if consteval}
statement that the frontend inside \texttt{nvcc} does not implement. The project
needs GCC 16 for \texttt{<mdspan>} and for OpenMP 5.2.

The resolution is a C ABI boundary. The \texttt{.cu} files are compiled by
\texttt{nvcc} driving GCC 14; everything else is compiled by GCC 16; and every
CUDA entry point is declared \texttt{extern "C"} taking plain pointers and
scalars. No libstdc++ type crosses, so the two compilers never have to agree on
a C++ ABI, only on the platform C ABI, which they do by definition. This also
satisfies the rule that backend files never leak their model's types.

The device solver is not a \texttt{Backend} implementation, and that is the right
call rather than an omission. The entire point of running on a GPU is that state
stays in device memory for the whole solve; forcing it behind
\texttt{parallel\_for} would mean a host callback per chunk, which would measure
PCIe latency and nothing else. The grid crosses once at the start and the answer
once at the end, and the transfer time is reported separately so that a reader
whose use case is a single small solve can add it back.

Natural ordering Gauss Seidel is absent from the device on purpose. It is
sequentially dependent and has no parallelism to offer a wide machine. That is a
result, not a gap in the implementation.

\section{Thread pinning under a hypervisor}

The library offers five binding policies, but two of them depend on being able
to tell a performance core from an efficiency core, and on this machine that
cannot be done from inside the guest.

The host is an i7-14700K with eight performance cores of two threads each and
twelve efficiency cores. The guest reports fourteen uniform cores of two
threads, identical 2048K L2 on every logical processor, \texttt{cpu\_capacity} of
1024 everywhere, no hybrid flag in \texttt{/proc/cpuinfo}, and no
\texttt{cpufreq} directory at all. Separately, \texttt{sched\_setaffinity} binds
a thread to a guest virtual processor, and the hypervisor remains free to place
that processor on any host core: a binding is a hint about the guest, not a
guarantee about silicon.

The library therefore measures instead of asking. A probe times an identical
compute kernel on each logical processor and reports a two group split only when
the throughputs separate by a margin that dominates the within group spread;
otherwise it says so, and the core type policies decline to bind. The knee that
Objective 3 asks for is taken from the aggregate scaling curve, which depends
only on how many cores are engaged and not on which, and is therefore valid
under virtualisation. Section \ref{sec:knee} reports both outcomes.
